\begin{textsample}{POS Dim 1 – vem\_ed\_10 – Score 14.00 – vem\_ed\_10/t043.txt}  \label{ex:f1_pos_119}
Jornada de PCIs: \textbf{compartilhando} experiências “We are SOME of the champions” reuniu relatos apresentados por professores com experiência \textbf{consolidada} em diversos tipos de PCIs.
Camila Kami, professora de inglês, tem experiências com PCIs desde 2019, nas Fatecs Bauru e Garça.
Ressaltou a \textbf{importância} de ferramentas tecnológicas user-friendly, como o Padlet.
Na atividade de ice-breaker, sugeriu um piquenique virtual.
Recomendou reuniões semanais com o professor parceiro, conversas com os alunos sobre \textbf{questões} interculturais e rubricas com critérios de avaliação.
Elizabeth Herrera Colorado, professora de espanhol, relatou um projeto sobre \textbf{interculturalidade} da Fatec Itaquaquecetuba com a UDEM (México). “Paciência, \textbf{flexibilidade} e acompanhamento permanente \textbf{são} nossos aliados”, resumiu.
Ela recomenda ainda entrar em sintonia com os parceiros internacionais e preparar os alunos para as apresentações \textbf{síncronas}.
Ricardo Pompeu comentou o PCI \textbf{que} as Fatecs Americana e Sumaré desenvolveram com Symbiosis sobre satisfação no trabalho, envolvendo 100 alunos (50 de \textbf{cada} país).
Segundo Pompeu, preparação prévia, tanto dos professores \textbf{quanto} dos alunos, \textbf{foi} \textbf{essencial} para o sucesso. “O PCI vai muito além da ementa da \textbf{disciplina}”, frisou.
Neusa Haruka \textbf{contou} sobre o PCI com Tianjin Normal University (China), sobre \textbf{aprendizado} do inglês como língua estrangeira.
O desafio principal \textbf{é} o fuso \textbf{horário} de 11 horas, o \textbf{que} \textbf{permitia} reuniões \textbf{síncronas} apenas nos fins de semana.
Mesmo assim, os alunos participavam.
A atividade intercultural consiste em pesquisas sobre a cultura do outro país.
Esse projeto tem três anos e meio e envolve 8 Fatecs.

% matched lemmas: aprendizado, cada, compartilhar, consolidar, contar, disciplina, essencial, flexibilidade, horário, importância, interculturalidade, permitir, quanto, que, questão, ser, síncrono
\end{textsample}
